\section{运动与外观分离：Two-Stream 网络和光流}

\subsection{为什么要把 motion 和 appearance 拆开}

视频里的信息并不总是适合混在一起处理。很多动作识别任务里，外观和运动各自携带不同类型的线索：

\begin{itemize}
    \item 外观分支更像图像分类器，负责识别人、物体和背景。
    \item 运动分支更像动作分类器，负责识别帧与帧之间的变化。
\end{itemize}

把这两种信号拆开，往往比强行一锅炖更有效。这就是 two-stream 的核心直觉。

\begin{figure}[H]
\centering
\includegraphics[width=0.95\textwidth]{slides-images/slide-043.jpg}
\caption{运动信息本身就足以让人理解动作：人类对 biological motion 的感知很强\protect\footnotemark}
\end{figure}
\footnotetext{Slides 第44页。}

\subsection{光流到底在测什么}

Two-stream 最经典的运动表示是 optical flow。它不是直接输入 RGB，而是输入相邻帧之间每个像素的位移估计。更形式化地说，光流是一个位移场

\[
F(x, y) = (d_x, d_y),
\]

其中 $d_x$ 和 $d_y$ 分别表示水平和垂直方向的位移。直觉上，光流告诉我们“这个像素在下一帧会往哪里动”。在经典亮度恒常假设下，可以写成

\[
I_{t+1}(x + d_x, y + d_y) \approx I_t(x, y).
\]

\begin{knowledgebox}{为什么光流有用}
\begin{itemize}
    \item 它把“运动”显式化了。
    \item 它让模型少做一层“自己猜运动”的工作。
    \item 它对动作识别特别有效，因为很多动作的核心就是动态变化，而不是静态外观。
\end{itemize}
\end{knowledgebox}

\begin{figure}[H]
\centering
\includegraphics[width=0.95\textwidth]{slides-images/slide-044.jpg}
\caption{光流定义：从帧 $t$ 到帧 $t+1$ 的像素位移场\protect\footnotemark}
\end{figure}
\footnotetext{Slides 第45页。}

\begin{figure}[H]
\centering
\includegraphics[width=0.95\textwidth]{slides-images/slide-045.jpg}
\caption{光流在两个方向上的分量可以分别可视化成 horizontal flow 和 vertical flow\protect\footnotemark}
\end{figure}
\footnotetext{Slides 第46页。}

\subsection{Two-stream 的架构}

Two-stream 的结构很直接：

\begin{enumerate}
    \item \textbf{Spatial stream}：输入单帧 RGB 图像，学外观。
    \item \textbf{Temporal stream}：输入多帧光流堆叠，学运动。
    \item \textbf{Fusion}：把两个分支的预测平均、加权或用分类器再融合。
\end{enumerate}

\begin{figure}[H]
\centering
\includegraphics[width=0.95\textwidth]{slides-images/slide-046.jpg}
\caption{Two-stream 的输入形式：一个分支看单帧图像，另一个分支看光流堆叠\protect\footnotemark}
\end{figure}
\footnotetext{Slides 第47页。}

\begin{figure}[H]
\centering
\includegraphics[width=0.95\textwidth]{slides-images/slide-047.jpg}
\caption{Two-stream 在 UCF-101 上的结果：运动分支、外观分支和融合后效果有明显差别\protect\footnotemark}
\end{figure}
\footnotetext{Slides 第48页。}

\begin{importantbox}{Two-stream 的经验现象}
课程中强调了一个反直觉但重要的现象：只看 motion 的 temporal stream 往往已经很强，有时甚至比单纯的 spatial stream 更强。这说明动作任务里，运动信号可能比静态外观更关键。
\end{importantbox}

\subsection{Two-stream 的局限}

Two-stream 解决了“外观和运动混在一起不容易学”的问题，但它也有明显代价：

\begin{itemize}
    \item 光流本身要预计算，工程成本高；
    \item 模型是两个分支，整体更复杂；
    \item 两个分支的信息融合还比较粗。
\end{itemize}

这也是后续 3D CNN、I3D、Non-local block 和 Transformer 风格方法出现的背景。

\begin{warningbox}{Two-stream 的经验教训}
显式运动表示很强，但它把前处理和模型切成了两部分。随着任务变大、视频变长、数据变多，大家越来越希望模型内部自己学会时空表征。
\end{warningbox}

